\section{The magnetic recursion in two squares}\label{sec:magnetic}

The bulk trajectory controls activities far from the boundary and the
source.  We now extend that control to the free walls and the magnetic
mark.  The purpose of the extension is to compare two squares using
exactly the same scalar normalization at the mark.  The insertion at
entry can have a polynomially large norm; its residual part will
contract, while its extracted scalar will cancel in the comparison.

Fix a large integer $j$, put $M=L^j$, and consider the free dual
squares of sides $Q=M$ and $Q=LM$, with their common center at the
mark.  Choose the $J$-history of Proposition~\ref{prop:bulk-window},
where $J$ is a power of two and $J\le j<2J$.  Its entry scale is
$m=L^i=qRs$, with $s\asymp_L J^3$ and $i=O(\log J)$.
In both squares stop at the same scale
\begin{equation}\label{eq:mag-common-stop}
 v=L^h,\qquad
 j-h=\left\lceil\log_L(j^{1/10})\right\rceil,
 \qquad D=\frac{M}{v}.
\end{equation}
Thus $j^{1/10}\le D<Lj^{1/10}$ and $M/(Ps)\to\infty$.
In particular the entry footprint and every fixed number of blocks
about the mark fit strictly inside both squares at all scales used.

We use the blocks, compatibility relation, analytic activity norms,
and covariance shells of the preceding sections.  Write
$\mathscr Z_k(K;\phi)$ for the subset partition functional at scale
$L^k$: it sums products of the activities $K(X;\phi)$ over compatible
collections of nonempty connected supports, including the empty
collection.  Thus $\mathscr Z_k$ is the functional
$\mathcal Z_{L^k}$ of Section~\ref{sec:inputs}, indexed here by its
absolute scale.  For another activity family $I$, define
\begin{equation}\label{eq:mag-first-variation}
 \mathscr I_k(K,I;\phi)
 =\left.\frac{d}{d\varepsilon}
       \mathscr Z_k(K+\varepsilon I;\phi)\right|_{\varepsilon=0}.
\end{equation}
Each term on the right has exactly one $I$ factor; the original
exclusions between that factor and the background factors are
retained.  The entry representation of Proposition~\ref{prop:entry}
has this form in the pierced square: in its notation,
$\mathscr I_i(K,I;\phi)=\mathcal Z_m^{(1)}(K,I;\phi)$.

For every support $X$ at scale $L^k$, let $\operatorname{Rect}_k(X)$ be its
axis-parallel spanning rectangle, and let $\operatorname{Rect}_k^H(X)$ be
that rectangle enlarged by $HL^k$ in each coordinate direction.
In a square, call $X$ \emph{bulk-good} if $\operatorname{Rect}_k^H(X)$ lies
strictly inside the square, and a \emph{wall support} otherwise.
Independently, call $X$ \emph{shift-good} if
$\operatorname{Rect}_k^H(X)$ misses the mark, and a \emph{defect support}
otherwise.  These predicates apply to supports of every size; they
specify both the full wall and defect norms below and the terminal
support decomposition.

The constant $H$ is fixed before $L$, independently of the entry
box parameter $P$.  It is large enough for the testing and assignment
radii and for the geometric sum of predecessor-dependence radii
used in copying.  The same definition at an output scale uses
$H$ \emph{coarse} blocks.  It does not require a bulk-good fine
input to be clear of its entire shell, whose range can be of order
$L$ fine blocks.  The entry convention for unused background supports
inside the compulsory footprint allows $H$ to remain independent
of $P$.

Figure~\ref{fig:common-core} shows the geometry of the common stop.
The local computations determining the marked scalars take place in
the same interior region; the wall contributions remain in the
terminal activities.
\begin{figure}[tb]
\centering
\begin{tikzpicture}[x=0.8cm,y=0.8cm,>=Latex,font=\small]
 \fill[gray!9] (-3.6,-3.6) rectangle (3.6,3.6);
 \fill[white] (-3.3,-3.3) rectangle (3.3,3.3);
 \draw[thick] (-3.6,-3.6) rectangle (3.6,3.6);
 \draw[thick] (-2.4,-2.4) rectangle (2.4,2.4);
 \draw[densely dashed] (-2.12,-2.12) rectangle (2.12,2.12);
 \fill[blue!8] (-0.58,-0.58) rectangle (0.58,0.58);
 \draw[blue!55!black,thick] (-0.58,-0.58) rectangle (0.58,0.58);
 \fill (0,0) circle (2pt);
 \node[above right] at (0,0) {$2\pi$};
 \draw[->] (0.58,0.58)--(4.1,1.8);
 \node[right,align=left,text width=4cm] at (4.1,1.8)
  {Common interior data\\within distance $O(v)$\\determine extracted scalars.};
 \draw[->] (2.38,-1.5)--(4.1,-0.7);
 \node[right,align=left,text width=4cm] at (4.1,-0.7)
  {Wall terms remain in\\the terminal comparison.};
 \draw[<->] (-2.4,-2.75)--(2.4,-2.75)
  node[midway,fill=white] {$M$};
 \draw[<->] (-3.6,-4)--(3.6,-4)
  node[midway,fill=white] {$LM$};
 \node[above] at (0,3.75) {$v=L^h,\qquad D=M/v\longrightarrow\infty$};
\end{tikzpicture}
\caption{Two concentric dual squares are stopped at the same block side
\(v\). The marked scalar computations use identical interior data,
whereas the remaining backgrounds are compared by terminal integration.
The common interior region has radius \(O(v)\), much smaller than \(M\).
The squares and collars are schematic and not drawn to scale.}
\label{fig:common-core}
\end{figure}


Let $K^0_{k,Q}$ and $K^l_{k,Q}$ denote the ordinary and pierced
background activities.  On a shift-good support the local lift
$\lambda_Q$ satisfies
$\partial\lambda_Q=\eta_l/\alpha$ and is defined up to a constant
multiple of the field period.  On interior supports meeting the
mark, the comparison with the pierced plane calculation instead
uses the smooth field
\[
 \gamma_Q=(r_Q-r_\infty)/\alpha.
\]
The lift estimates and entry copying established earlier apply in
these field units.

\begin{proposition}[Common magnetic normalization]\label{prop:magnetic}
For all sufficiently large $j$, the exact maps from entry scale $i$
to stop $h$ can be chosen with the following properties.
\begin{enumerate}
\item At the stop, in the fixed norm regimes for the ordinary and
defect activities,
\begin{equation}\label{eq:mag-small-backgrounds}
 \max_{Q\in\{M,LM\}}
 \bigl(\norm{K^0_{h,Q}}+\norm{K^l_{h,Q}}\bigr)\le \frac{C}{j}.
\end{equation}
The bulk copies satisfy the sharper conclusion of
Proposition~\ref{prop:bulk-window}.  On shift-good supports, pierced
activities are exactly ordinary activities evaluated at the local
lift.  Interior pierced activities copy the plane calculation at
the translate $\gamma_Q$, whenever their required clearance holds.

\item Let $a^\nu_{k,Q,b}$ be the designated scalar coefficient at
coarse block $b$ in step $k$, for $\nu=0,l$, and set
\begin{equation}\label{eq:mag-background-products}
 P^\nu_Q=\prod_{k=i}^{h-1}\prod_b(1+a^\nu_{k,Q,b}).
\end{equation}
All factors are real and positive.  The ratios agree exactly:
\begin{equation}\label{eq:mag-common-background-ratio}
 \frac{P^l_M}{P^0_M}=\frac{P^l_{LM}}{P^0_{LM}}=:\rho_j.
\end{equation}
For every fixed $\epsilon_2>0$, the initial smallness choices can
be made so that $|\log\rho_j|\le\epsilon_2j$ for all sufficiently
large $j$.

\item There are a real scalar $\sigma_Q$ and residual insertion
activities $I_{h,Q}$ such that integration of the entry insertion
over all shells below $v$ gives the exact field identity
\begin{equation}\label{eq:mag-insertion-identity}
 \E_{i:h}\mathscr I_i(K^l_{i,Q},I_{i,Q};\phi)
 =P^l_Q\left\{
    \sigma_Q\mathscr Z_h(K^l_{h,Q};\phi)
       +\mathscr I_h(K^l_{h,Q},I_{h,Q};\phi)\right\}.
\end{equation}
Here $\E_{i:h}$ denotes successive shell integration, with the
remaining field $\phi$ held fixed.  The scalar is common to the
two squares,
\begin{equation}\label{eq:mag-common-insertion-scalar}
 \sigma_M=\sigma_{LM}=: \sigma_j,
\end{equation}
and every support of $I_{h,Q}$ contains the designated mark block.
Moreover,
\begin{equation}\label{eq:mag-insertion-bound}
 \norm{I_{h,Q}}\le C_Rs^{C_R}L^{-2(h-i)/3},
 \qquad Q\in\{M,LM\}.
\end{equation}
\end{enumerate}
The products in \eqref{eq:mag-background-products} contain only
the RG extraction factors.  They do not include the common entry
scalar or the harmonic-energy factor of the Gaussian sector law.
\end{proposition}

We prove the proposition in four steps.  The localizations and
placement estimates are those of the exact partition map in
\cite[Section 4]{B}.  The wall and lifted-copy
constructions follow the method of its Propositions
5.5 and 5.7.  Their
low-temperature conclusions are not inputs here: the marginal
charge estimate, critical forcing, and insertion contraction needed
for the present proposition are established below.

\subsection{The wall estimate at the marginal frequency}

A small support means one in the small linear-input class of the
partition map.  Its fine-block size is bounded by the fixed
threshold $n_0$, and its coarse closure has index diameter at most
one.  In an ordinary square use the plane localization, including
the fundamental harmonic localization, on small bulk-good inputs.
On small wall inputs extract only the neutral constant.  Since a
small support has bounded diameter, the latter inputs lie within
a fixed number of fine-block layers of a wall.  A small bulk-good
input still uses the plane prescription if its shell convolution
can see the wall; its contribution to a wall output is treated as
forcing.  In contrast, the $H$ coarse-block clearance of a bulk-good
output ensures that every predecessor needed for copying lies in
the interior with its required shell and numerical dependence.

Ordinary activities are even in the field.  After neutral constant
subtraction, the localization lemma of
\cite[Lemma 4.5]{B} gives a factor $CL^{-2}$
for a small wall input.  There are $O_H(L)$ fine placements at a
prescribed coarse block along a wall, and only $O_H(1)$ at a
corner.  It remains to check the charged part at
$\alpha^2=8\pi$.

Let $\Gamma$ be the shell covariance from side $m'$ to side $Lm'$,
and let $x$ lie on a small wall support.  The binary kernel
estimates, with even reflection at a free wall, give
\begin{equation}\label{eq:mag-wall-shell}
 \Gamma(x,x)\ge\frac{\log L}{2\pi}-C,
 \qquad
 \norm{\Gamma(\,\cdot\,,x)-\Gamma(x,x)}_{m',X}\le C.
\end{equation}
Both constants are independent of $L$ after the geometric radii
are fixed.  For the diagonal estimate, the direct image supplies
the plane value.  An image contribution vanishes at binary scales
smaller than a fixed multiple of its separation, by finite range.
At scales larger than a sufficiently large fixed multiple of that
separation it is nonnegative: the positive plane binary diagonal
and the derivative estimate give this assertion.  Only boundedly
many intermediate binary scales can contribute negatively, and
each is bounded.  At a corner the same argument applies to the
finite set of reflected images.  Other walls are beyond the shell
range in the applications where this description is used.
Positive-order binary estimates summed over the shell prove the
second bound in \eqref{eq:mag-wall-shell}.

Apply the charge-contour estimate of
\cite[Equations (4.26) and
(8.14)]{B}, with shift amplitude
$a=\alpha$.  For nonzero integer charge $q$, its exponential gain
is bounded by
\[
 C\exp\{-\alpha^2(|q|-1/2)\Gamma(x,x)+h_0\alpha|q|\}.
\]
Using \eqref{eq:mag-wall-shell}, the fundamental charges have gain
$CL^{-2}$, and the higher charges form a summable series with
stronger powers of $L^{-1}$.  The analytic radius $h_0$ is chosen
before $L$ to allow the fixed nonconstant imaginary shift.  The
observation regulator is admissible in this estimate by
\cite[Lemma 9.2]{B} and its free-square version
verified at entry.  Thus the charged wall terms have the same
$CL^{-2}$ bound required for the placement count.

Let $p_k$ be the full ordinary plane activity norm, and let $w_k$
be the full norm restricted to ordinary wall supports as defined
above, without a small-size restriction.  Large linear inputs
have arbitrarily strong weight contraction after the activity
weight is chosen sufficiently large.  The analytic nonlinear
bound of the partition map and the preceding small-input estimates
give
\begin{equation}\label{eq:mag-wall-recursion}
 w_{k+1}\le\theta_1w_k+C_Lp_k+C_L(w_k+p_k)^2,
 \qquad \theta_1<1.
\end{equation}
Here $C_Lp_k$ includes plane inputs and assigned pieces that reach
a wall output.  Their field evaluations occur on their original
fine blocks with the prescribed collars, so no wall cancellation
is needed to bound them.  All norm regimes, including their
regulator reserves, are fixed once; the wall contraction does not
consume another reserve at each step.

\subsection{Defect localization and exact copying}

On a small shift-good input in the pierced square, use exactly
the localization pieces of its ordinary counterpart, evaluated at
the lifted field.  Designate exactly the same scalar for
extraction.  In particular a translated quadratic may have a
nonzero value at zero field, but that value remains part of the
quadratic piece; it is not extracted as an additional scalar.

For a small defect input $F_Q$, let
\[
 (F_Q)_0(\phi)=\frac{\alpha}{2\pi}
       \int_0^{2\pi/\alpha}F_Q(\phi+t)\,dt
\]
be its neutral projection.  In the plane extract
$\E(F_\infty)_0(\zeta)$, where $\zeta$ is the current shell.
In the square extract instead
\begin{equation}\label{eq:mag-backtranslated-extraction}
 c_Q(F_Q)=\E(F_Q)_0(\zeta-\gamma_Q).
\end{equation}
The neutral projection is invariant under constant field shifts,
so the alignment constant in $\gamma_Q$ does not affect this
definition.  Every small defect input is interior: its bounded
fine-block diameter and its failure to be shift-good place it
within $O_H(L^k)$ of the mark, while $L^k\le v/L=M/(DL)$ for
every input to a step before the stop.

The changed subtraction has the same contraction as ordinary
neutral constant subtraction.  Indeed the latter gives $CL^{-1}$
without evenness.  On a small fine support at side $m'=L^k$, the
nonconstant part of $\gamma_Q$ has fine test norm at most
$Cm'/Q$.  By the analytic derivative bound with its fixed
regulator reserve,
\begin{equation}\label{eq:mag-extraction-change}
 \left|c_Q(F_Q)-\E(F_Q)_0(\zeta)\right|
 \le C\frac{m'}{Q}\norm{F_Q}
 \le \frac{C}{L}\norm{F_Q}.
\end{equation}
The last inequality holds uniformly through the chosen stop for
large $j$.  Charged small defect inputs have gain $CL^{-2}$ by the
same contour argument, now with the interior shell estimates.
At a prescribed output there are only $O_H(1)$ small fine defect
placements.  Consequently their total residual norm is at most
$CL^{-1}$ times their input norm.

Translations of shift-good inputs use the norm reserves fixed at
entry.  On the relevant fine collars,
$m'|\nabla\lambda_Q|$ and $(m')^2|\nabla^2\lambda_Q|$ are bounded.
The quadratic inequality therefore permits a fixed increase in
the ordinary regulator exponent, and a fixed increase in the
coefficient of the variational observation regulator.  Each costs
at most $C^{|X|}$ on a support $X$, absorbed by the reserved activity
weight.  A localized piece that is assigned to a coarse singleton
is still evaluated on the fine block where its lift is defined.
This applies to both gradient pieces and fundamental harmonics;
a deterministic real phase does not enlarge the analytic norm of
the latter.  Thus pieces reaching defect outputs have controlled
norm without evaluating a lift at the singular core.

Let $d_k$ be the full defect norm, in this fixed reserved regime.
The same multilinear partition-map estimates as for the wall,
with the large-input weight gain in the defect norm, give
\begin{equation}\label{eq:mag-defect-recursion}
 d_{k+1}\le\theta_1d_k+C_L(p_k+w_k)
             +C_L(d_k+w_k+p_k)^2,
 \qquad\theta_1<1.
\end{equation}
The forcing term includes the translated ordinary pieces just
described.  This proves the analytic estimate needed near the mark.
We next verify the numerical equality that will make its scalar
normalization disappear between the two squares.

Suppose an output support has the stated clearance.  Every fine
input, assigned singleton piece, and scalar normalization entering
that output then fits in its enlarged neighborhood.  On a
shift-good output they share a common lift.  Products and Gaussian
convolution commute with this deterministic translation, and the
identical localizations preserve the translated-copy identity.
For an interior defect output use instead
\begin{equation}\label{eq:mag-pierced-copy}
 F_Q(\phi)=F_\infty(\phi+\gamma_Q).
\end{equation}
The shell kernels agree there by finite range, and
\eqref{eq:mag-backtranslated-extraction} gives the exact equality
\[
 c_Q(F_Q)=\E(F_\infty)_0(\zeta).
\]
It follows inductively that both activities and designated
coefficients copy the corresponding plane calculation.

This induction includes numerical, as well as field, dependence.
In the exact map, predecessor dependence from relocation and
singleton normalizers extends by at most a fixed multiple of the
current scale.  Over earlier steps these distances sum to a
geometric series.  The link construction at entry records any
numerical dependence reaching a wall in its support, so the same
argument starts at entry.  Enlarging $H$ by a fixed amount pays
these causal radii.  A support that still reaches a wall cannot
become a small marked input before the stop: the distance from
the mark to the wall, in stop blocks, tends to infinity.  Large
inputs do not supply scalar extractions.  Thus neither a long
entry link nor a large insertion support can introduce a hidden
wall dependence into the local constants at the mark.

The only differences between pierced and ordinary designated
constants arise from small defect inputs, so they occur at
$O_H(1)$ blocks per step.  At each of these blocks both the
pierced and ordinary constants are the same plane constants in
the two squares.  All other factors cancel in the
pierced-to-ordinary product within each square.  This proves
\eqref{eq:mag-common-background-ratio} exactly.

\subsection{Scalar extraction with one insertion}

For clarity, we record the scalar algebra before estimating the
insertion.  Write one exact background map as
\begin{equation}\label{eq:mag-one-map}
 \E\mathscr Z_k(K;\phi+\zeta)
   =p_k(K)\mathscr Z_{k+1}(\mathcal F_k(K);\phi),
 \qquad
 p_k(K)=\prod_b(1+a_{k,b}(K)).
\end{equation}
Here the $a_{k,b}$ are the designated coefficients of small linear
inputs.  Extend the same exact splitting algebra to coefficients in
$\R[\varepsilon]/(\varepsilon^2)$.  Its degree-zero terms use the
background localization; its degree-one terms use the marked
localization and allocation specified below.  Write
\[
 a_{k,b}(\varepsilon)=a_{k,b}+\varepsilon a'_{k,b},
 \qquad K_+=\mathcal F_k(K),\qquad
 K_+(\varepsilon)=K_++\varepsilon\mathcal T_k(K)I.
\]
Thus $\mathcal T_k(K)$ is the linear map in the insertion slot of
this extended construction.  It need not be the derivative of the
unmarked map $\mathcal F_k$, since the two slots have different
localization prescriptions.  Put
\[
 I_+=\mathcal T_k(K)I,
 \qquad
 \beta_k=\sum_b\frac{a'_{k,b}}{1+a_{k,b}}.
\]
The coefficient of $\varepsilon$ in the exact identity is
\begin{equation}\label{eq:mag-differentiated-map}
 \E\mathscr I_k(K,I;\phi+\zeta)
 =p_k(K)\left\{\beta_k\mathscr Z_{k+1}(K_+;\phi)
                    +\mathscr I_{k+1}(K_+,I_+;\phi)\right\}.
\end{equation}
This is differentiation of the partition identity with the stated
degree-dependent splittings.  No finite nonzero perturbation by the
possibly large $I$ is required.

If an earlier step has already produced a scalar $\sigma$, linearity
applied to $\sigma\mathscr Z_k(K)+\mathscr I_k(K,I)$ gives
\begin{equation}\label{eq:mag-scalar-recursion}
 \sigma_+=\sigma+\beta_k,
 \qquad I_+=\mathcal T_k(K)I.
\end{equation}
In particular the earlier scalar is kept outside the later map.
It is not put back into the activity as $\sigma K$.
This elementary identity is important because $\sigma$ need not
be small, whereas the operator controlling $I_+$ only sees the
small background $K$.

We also specify how the mark is preserved in the exact algebra.
Use \eqref{eq:mag-backtranslated-extraction} for small insertion
inputs, and assign their extracted constant to the coarse block
containing the designated fine mark block.  This assignment lies
in their coarse closure, as required by the partition map.  A
residual input keeps its closure label, including the mark.
In the scalar normalization step the decoration at block $b$ is
\[
 d_b=(1+a_b)^{-1}-1,
 \qquad d'_b=-\frac{a'_b}{(1+a_b)^2}.
\]
Only a small marked input can give $a'_b\ne0$, and its designated
block is the mark block.  Hence a differentiated decoration is
also supported at the mark.  Every term in a differentiated
product contains either an insertion input or such a decoration;
its final connected support therefore contains the mark.  Original
exclusions remain attached to relocated pieces, exactly as in
\cite[Proposition 4.8]{B}.  The differentiation
does not turn a compatible collection into an unrestricted product.

At zero background, the complete residual insertion map after
scalar extraction is the subtracted small linear map together
with the unchanged large linear map.  The small-input bound just
proved is $C/L$, with boundedly many marked placements.  By
increasing the large-support weight after fixing $L$, its remaining
linear contribution is as small as desired.  The multilinear bounds
for the exact partition map give
$\norm{\mathcal T_k(K)-\mathcal T_k(0)}\le C_L\norm{K}$ on the
fixed small background ball, and hence
\begin{equation}\label{eq:mag-insertion-operator}
 \norm{I_+}
 \le\left(\frac{C}{L}+\delta_A
                       +C_L\norm{K}\right)\norm{I},
\end{equation}
where $\delta_A$ tends to zero as the weight is increased.
For this degree-one map the background and insertion may have
different designated localizations: the exact splitting algebra
is linear in each input, and its multilinear estimates use the
same bounds on these designated pieces.  Thus
\eqref{eq:mag-insertion-operator} applies to the prescribed marked
splitting as well.  It imposes no smallness assumption on
$\norm{I}$ itself.

Choose $L$ so that $C/L<\tfrac13L^{-2/3}$, the weight so that
$\delta_A<\tfrac13L^{-2/3}$, and then the background tube so that
$C_L\norm{K}<\tfrac13L^{-2/3}$ at every step.  Iteration and the
entry insertion bound give \eqref{eq:mag-insertion-bound}.

Finally every summand in $\beta_k$ depends only on a small marked
input and the background normalizer at its assigned block.
The copying argument in the preceding subsection makes both its
numerator and denominator identical in the two squares.  This
also holds for the insertion activities by induction through
\eqref{eq:mag-differentiated-map}.  At entry the compulsory footprint
fits with clearance in both squares; a support too small to
contain it has insertion coefficient zero in both.  Starting with
$\sigma_i=0$, Equation~\eqref{eq:mag-scalar-recursion} therefore
gives $\sigma_M=\sigma_{LM}$.  This equality does not require a
bound or a sign for the accumulated scalar.

\subsection{Completion of the background bounds}

\begin{proof}[Proof of Proposition~\ref{prop:magnetic}]
The constructions above give the exact identities and the copying
assertions.  We complete the uniform estimates and the parameter
choices.  Fix the kernel and assignment radii, the clearance $H$,
and the finite list of norm regimes with their translation and
integration reserves.  Choose $L$ and then the activity weights
as in the local estimates.  The constants in the small-support
bounds before placement counting are independent of $L$.
After these choices take the bulk diagnostic tube sufficiently
small.  Proposition~\ref{prop:entry} supplies arbitrarily small
entry backgrounds in all of these fixed regimes, and the bulk
analysis keeps the plane activities in the prescribed tube.

The two recursions \eqref{eq:mag-wall-recursion} and
\eqref{eq:mag-defect-recursion} now keep the wall and defect
backgrounds small throughout the history.  More explicitly, by
decreasing the tube the quadratic terms can be absorbed into a
slightly larger contraction $\theta_2<1$ and a constant multiple
of their forcing terms.  Iteration first for $w_k$ and then for
$d_k$ proves the assertion.  On the last positive-fraction window,
Proposition~\ref{prop:bulk-window} gives $p_k\le C/j$.
The geometric convolutions in the two recursions therefore give
$w_h+d_h\le C/j$; the earlier part of each convolution is
exponentially small in $j$.  This proves
\eqref{eq:mag-small-backgrounds}, including its uniformity for
$Q=M$ and $Q=LM$.

All coefficients and localizations are real.  Their scalar
coefficients are bounded by a fixed map constant times the small
background tube, so $|a^\nu_{k,Q,b}|<1/2$ and every extracted
factor is positive.  Only boundedly many factors per step differ
between the pierced and ordinary maps.  Since
\[
 |\log(1+a)-\log(1+a')|\le2|a-a'|
 \qquad (|a|,|a'|<1/2),
\]
their cumulative ratio has logarithm at most $\epsilon_2j$ in
absolute value after the tube is decreased in terms of the
prescribed $\epsilon_2$.  This decrease is made after the fixed map
constants and remains compatible with the preceding choices.

The same smallness gives the insertion operator bound chosen
after \eqref{eq:mag-insertion-operator}.  Iterating
\eqref{eq:mag-differentiated-map}, keeping the scalar outside each
subsequent map as in \eqref{eq:mag-scalar-recursion}, proves
\eqref{eq:mag-insertion-identity}.  The equality of its scalar
between the two squares and the norm bound have already been
established.  This proves all parts of the proposition.
\end{proof}
